% Focused framework presentation.
\section{Supplementary derivations}

\begin{frame}[t]{Appendix A1: algebra operations}
\hypertarget{app:algebras}{}
\begin{columns}[T,onlytextwidth]
  \begin{column}{.48\textwidth}
    \heading{The algebra $\operatorname{End}(V)$}
    \begin{equation*}
      \begin{aligned}
        (A+B)x&=Ax+Bx,\\
        (aA)x&=a(Ax),\\
        (AB)x&=A(Bx).
      \end{aligned}
    \end{equation*}
    Composition is associative, with unit $I_V$;
    $AB$ and $BA$ can differ when $d\geq2$.
    \par\medskip
    \flowdiagram{composition}
  \end{column}
  \begin{column}{.47\textwidth}
    \heading{The observable algebra $\mathcal A=B_b(V)$}
    \begin{equation*}
      \begin{aligned}
        (af+bg)(x)&=af(x)+bg(x),\\
        (fg)(x)&=f(x)g(x),\\
        1(x)&=1.
      \end{aligned}
    \end{equation*}
    Bounded Borel functions form a commutative unital algebra.
    \par\medskip
    \flowdiagram[.80\linewidth]{observables}
  \end{column}
\end{columns}
\slidecite{Manuscript \S\,1.1. \hyperlink{main:states}{Return to state space and algebra.}}
\end{frame}

\begin{frame}[t]{Appendix A2: existence and uniqueness}
\hypertarget{app:existence}{}
\begin{columns}[T,onlytextwidth]
  \begin{column}{.53\textwidth}
    \heading{Bounds on a compact time window}
    For $v_t(y)=A_ty+B_t$ and compact $J\subset I$, let
    \begin{equation*}
      L_J=\sup_{t\in J}\|A_t\|,\qquad
      C_J=\sup_{t\in J}\|B_t\|.
    \end{equation*}
    Both constants are finite, and
    \begin{equation*}
      \begin{aligned}
        \|v_t(y)-v_t(z)\|&\leq L_J\|y-z\|,\\
        \|v_t(y)\|&\leq C_J+L_J\|y\|.
      \end{aligned}
    \end{equation*}
  \end{column}
  \begin{column}{.42\textwidth}
    \heading{Continuation argument}
    \begin{enumerate}\setlength{\itemsep}{7pt}
      \item Local Lipschitz continuity gives a unique local solution.
      \item Linear growth prevents finite-time escape on $J$.
      \item Uniqueness joins the solutions on overlapping windows.
    \end{enumerate}
    The trajectory is therefore uniquely defined on all of $I$.
  \end{column}
\end{columns}
\slidecite{Manuscript Example 1.10 and Appendix A.1. \hyperlink{main:ivp}{Return to the initial-value problem.}}
\end{frame}

\begin{frame}[t]{Appendix A3: evaluation and evolution}
\hypertarget{app:evaluation}{}
\begin{columns}[T,onlytextwidth]
  \begin{column}{.52\textwidth}
    \heading{Homogeneous trajectories}
    \begin{equation*}
      \mathcal H=\{\eta\in C^1(I,V):\dot\eta_t=A_t\eta_t\}.
    \end{equation*}
    Evaluation $E_s(\eta)=\eta_s$ is linear.
    \begin{itemize}\setlength{\itemsep}{4pt}
      \item Existence makes it onto.
      \item Uniqueness makes it one-to-one.
    \end{itemize}
    \begin{equation*}
      S_s=E_s^{-1},\qquad \mathcal H\cong V,
      \qquad\dim\mathcal H=d.
    \end{equation*}
  \end{column}
  \begin{column}{.43\textwidth}
    \heading{Inverse evaluation}
    \flowdiagram{evaluation}
    \begin{equation*}
      E_sS_s=I_V,\qquad S_sE_s=I_{\mathcal H}.
    \end{equation*}
  \end{column}
\end{columns}
\heading{Composition of the homogeneous evolution}
\begin{equation*}
  U(t,r)U(r,s)=E_tS_rE_rS_s=E_tS_s=U(t,s).
\end{equation*}
\slidecite{Manuscript Example 1.10. \hyperlink{main:homogeneous}{Return to homogeneous evolution.}}
\end{frame}

\begin{frame}[t]{Appendix A4: the affine solution family}
\hypertarget{app:affine}{}
\begin{columns}[T,onlytextwidth]
  \begin{column}{.52\textwidth}
    \heading{A particular solution shifts the kernel}
    Suppose $\mathsf D_A\gamma_{\mathrm p}=B$. Then
    \begin{equation*}
      \begin{aligned}
        \mathsf D_A\gamma=B
          &\iff\mathsf D_A(\gamma-\gamma_{\mathrm p})=0\\
          &\iff\gamma-\gamma_{\mathrm p}\in\mathcal H.
      \end{aligned}
    \end{equation*}
    Hence
    \begin{equation*}
      \mathcal S_B=\gamma_{\mathrm p}+\mathcal H.
    \end{equation*}
  \end{column}
  \begin{column}{.43\textwidth}
    \heading{The initial condition selects one path}
    Write $\gamma=\gamma_{\mathrm p}+\eta$. Then
    \begin{equation*}
      E_s\eta=x-\gamma_{\mathrm p}(s).
    \end{equation*}
    Since $E_s$ is an isomorphism,
    \begin{equation*}
      \gamma=\gamma_{\mathrm p}
        +S_s\bigl(x-\gamma_{\mathrm p}(s)\bigr).
    \end{equation*}
    Taking $\gamma_{\mathrm p}(s)=0$ gives
    \begin{equation*}
      \gamma=\gamma_{\mathrm p}+S_sx.
    \end{equation*}
  \end{column}
\end{columns}
\slidecite{Manuscript Example 1.10. \hyperlink{main:affine}{Return to the affine solution.}}
\end{frame}

\begin{frame}[t]{Appendix A5: variation of constants}
\hypertarget{app:variation}{}
\begin{columns}[T,onlytextwidth]
  \begin{column}{.48\textwidth}
    \heading{1. Vary the homogeneous coefficient}
    \begin{equation*}
      \gamma_{\mathrm p}(t)=U(t,s)c(t),\qquad c(s)=0.
    \end{equation*}
    \heading{2. Apply the trajectory operator}
    \begin{equation*}
      \begin{aligned}
        \dot\gamma_{\mathrm p}-A_t\gamma_{\mathrm p}
          &=U(t,s)\dot c(t),\\
        \dot c(t)&=U(s,t)B_t.
      \end{aligned}
    \end{equation*}
  \end{column}
  \begin{column}{.47\textwidth}
    \heading{3. Integrate and compose}
    \begin{equation*}
      c(t)=\int_s^t U(s,r)B_r\,\mathrm dr.
    \end{equation*}
    Composition gives $U(t,s)U(s,r)=U(t,r)$, hence
    \begin{equation*}
      \begin{aligned}
        \gamma_{\mathrm p}(t)
          &=\int_s^t U(t,r)B_r\,\mathrm dr.
      \end{aligned}
    \end{equation*}
  \end{column}
\end{columns}
\medskip
\begin{equation*}
  x_t=U(t,s)x+\int_s^t U(t,r)B_r\,\mathrm dr.
\end{equation*}
\slidecite{Manuscript Example 1.10. \hyperlink{main:affine}{Return to the solution formula.}}
\end{frame}

\begin{frame}[t]{Appendix A6: verification of the solution}
\hypertarget{app:verification}{}
\heading{Initial condition}
\begin{equation*}
  x_s=U(s,s)x+\int_s^s U(s,r)B_r\,\mathrm dr=x.
\end{equation*}
\heading{Differentiation}
Since $\partial_tU(t,r)=A_tU(t,r)$, Leibniz's rule gives
\begin{equation*}
  \begin{aligned}
    \dot x_t
      &=A_tU(t,s)x+U(t,t)B_t+\int_s^t A_tU(t,r)B_r\,\mathrm dr\\
      &=A_t\left(U(t,s)x+\int_s^t U(t,r)B_r\,\mathrm dr\right)+B_t\\
      &=A_tx_t+B_t.
  \end{aligned}
\end{equation*}
\heading{Uniqueness}
The difference of two solutions lies in $\mathcal H$ and has zero value at $s$.
Injectivity of $E_s$ makes that difference the zero trajectory.
\slidecite{Manuscript Example 1.10. \hyperlink{main:affine}{Return to the solution and its flow map.}}
\end{frame}

\begin{frame}[t]{Appendix A7: general evaluation and composition}
\hypertarget{app:general}{}
\begin{columns}[T,onlytextwidth]
  \begin{column}{.49\textwidth}
    \heading{Evaluation is a bijection}
    \begin{equation*}
      \mathcal S_v=\{\gamma\in C^1(I,V):\dot\gamma_t=v_t(\gamma_t)\}.
    \end{equation*}
    On the solution set, $E_s(\gamma)=\gamma_s$.
    Existence gives surjectivity; uniqueness gives injectivity.
    \begin{equation*}
      \Phi_{s,t}=E_t\circ E_s^{-1}.
    \end{equation*}
  \end{column}
  \begin{column}{.46\textwidth}
    \heading{Recover, then evaluate}
    \flowdiagram{general-evaluation}
  \end{column}
\end{columns}
\heading{Cancellation at an intermediate time}
\begin{equation*}
  \begin{aligned}
    \Phi_{r,t}\circ\Phi_{s,r}
      &=(E_t\circ E_r^{-1})\circ(E_r\circ E_s^{-1})\\
      &=E_t\circ E_s^{-1}=\Phi_{s,t}.
  \end{aligned}
\end{equation*}
\begin{equation*}
  \Phi_{s,s}=I_V,\qquad \Phi_{s,t}^{-1}=\Phi_{t,s}.
\end{equation*}
\slidecite{Manuscript Appendix A.2. \hyperlink{main:general}{Return to general flows}; \hyperlink{main:laws}{return to the flow laws}.}
\end{frame}

\begin{frame}[t]{Appendix A8: the determinant equation}
\hypertarget{app:jacobian}{}
\heading{Invertibility from the inverse flow}
\begin{equation*}
  D\Phi_{t,s}(\Phi_{s,t}(x))M_{s,t}(x)=I_d.
\end{equation*}
Thus $M=M_{s,t}(x)$ is invertible. Set $A_t=Dv_t(\Phi_{s,t}(x))$.
\heading{Jacobi's formula and cyclic invariance}
\begin{equation*}
  \begin{aligned}
    \partial_t\det M
      &=(\det M)\operatorname{tr}(M^{-1}\partial_tM)\\
      &=(\det M)\operatorname{tr}(M^{-1}A_tM)\\
      &=(\det M)\operatorname{tr}(A_t).
  \end{aligned}
\end{equation*}
Since $\operatorname{tr}(Dv_t)=\nabla\cdot v_t$ and $\det M_{s,s}=1$,
\begin{equation*}
  J_{s,t}(x)=\exp\!\left(\int_s^t\nabla\cdot v_r(\Phi_{s,r}(x))\dd r\right).
\end{equation*}
The determinant equation is scalar even when the matrices at different times fail to commute.
\slidecite{Manuscript \S\,1.2.1, variational matrices. \hyperlink{main:jacobian}{Return to local volume.}}
\end{frame}

\begin{frame}[t]{Appendix A9: derivatives of the Koopman operator}
\hypertarget{app:generator}{}
\heading{Terminal time}
For a smooth observable, the chain rule gives
\begin{equation*}
  \partial_t[f(\Phi_{s,t}(x))]
  =\nabla f(\Phi_{s,t}(x))\cdot v_t(\Phi_{s,t}(x))
  =(\mathcal K_{s,t}\mathcal L_tf)(x).
\end{equation*}
\heading{Initial time}
The composition law yields
\begin{equation*}
  \Phi_{s,t}=\Phi_{s+h,t}\circ\Phi_{s,s+h}.
\end{equation*}
Differentiate at $h=0$ to obtain
\begin{equation*}
  \partial_s\Phi_{s,t}(x)=-D_x\Phi_{s,t}(x)v_s(x).
\end{equation*}
Therefore,
\begin{equation*}
  \partial_s(\mathcal K_{s,t}f)(x)
  =-v_s(x)\cdot\nabla_x[f(\Phi_{s,t}(x))]
  =-\mathcal L_s(\mathcal K_{s,t}f)(x).
\end{equation*}
\slidecite{Manuscript \S\,1.2.2, generator identities. \hyperlink{main:generator}{Return to infinitesimal evolution.}}
\end{frame}

\begin{frame}[t]{Appendix A10: density change of variables}
\hypertarget{app:density}{}
For a Borel set $A$ and a $C^1$ diffeomorphism $\Phi=\Phi_{0,t}$,
\begin{equation*}
  \begin{aligned}
    \mu_t(A)&=\mu_0(\Phi^{-1}(A))\\
      &=\int_{\Phi^{-1}(A)}\rho_0(x_0)\dd x_0\\
      &=\int_A\rho_0(\Phi^{-1}(y))|\det D\Phi^{-1}(y)|\dd y.
  \end{aligned}
\end{equation*}
Thus $\rho_t(y)=\rho_0(\Phi^{-1}(y))|\det D\Phi^{-1}(y)|$ almost everywhere.
\heading{The inverse Jacobian}
\begin{equation*}
  D\Phi^{-1}(\Phi(x_0))D\Phi(x_0)=I_d,\qquad
  |\det D\Phi^{-1}(\Phi(x_0))|=J_{0,t}(x_0)^{-1}.
\end{equation*}
Substitution gives $\rho_t(\Phi(x_0))J_{0,t}(x_0)=\rho_0(x_0)$.
Only integrability of $\rho_0$ is needed for this formula. A pointwise classical
continuity equation requires the additional $C^1$ regularity of $\rho(t,x)$.
\slidecite{Manuscript \S\,1.2.2, Assessment 9.7--8. \hyperlink{main:density}{Return to density transport.}}
\end{frame}

\begin{frame}[t]{Appendix A11: the measure-valued weak equation}
\hypertarget{app:weak}{}
For $\psi\in C_c^\infty((0,T)\times\mathbb R^d)$, define
\begin{equation*}
  F(t)=\int\psi(t,x)\mu_t(\dd x)
  =\int\psi(t,\Phi_{0,t}(x_0))\mu_0(\dd x_0).
\end{equation*}
The chain rule along trajectories and local flux integrability give
\begin{equation*}
  F'(t)=\int[\partial_t\psi+v_t\cdot\nabla\psi](t,x)\mu_t(\dd x)
\end{equation*}
for almost every $t$. Integrating in time,
\begin{equation*}
  \int_0^T\int[\partial_t\psi+v_t\cdot\nabla\psi]\dd\mu_t\dd t=F(T)-F(0)=0.
\end{equation*}
The endpoint terms vanish because $\psi$ is supported away from $t=0,T$.
\heading{Distributional interpretation}
\begin{equation*}
  \langle\partial_t\mu+\nabla\cdot(v\mu),\psi\rangle
  =-\int_0^T\int[\partial_t\psi+v_t\cdot\nabla\psi]\dd\mu_t\dd t=0.
\end{equation*}
\slidecite{Manuscript \S\,1.2.2, weak continuity. \hyperlink{main:weak}{Return to probability transport without a density.}}
\end{frame}

\begin{frame}[t]{Appendix A12: symplecticity of the exact flow}
\hypertarget{app:symplectic}{}
Write $Q_t=\nabla^2H(z_t)=Q_t^\top$ and $\dot M_t=\mathbb J_nQ_tM_t$.
Since $\mathbb J_n^\top=-\mathbb J_n$ and $\mathbb J_n^2=-I$,
\begin{equation*}
  \begin{aligned}
    \frac{\dd}{\dd t}(M_t^\top\mathbb J_nM_t)
    &=M_t^\top[(\mathbb J_nQ_t)^\top\mathbb J_n
      +\mathbb J_n(\mathbb J_nQ_t)]M_t\\
    &=M_t^\top[Q_t(-\mathbb J_n)\mathbb J_n+\mathbb J_n^2Q_t]M_t\\
    &=M_t^\top(Q_t-Q_t)M_t=0.
  \end{aligned}
\end{equation*}
The initial condition gives $M_t^\top\mathbb J_nM_t=\mathbb J_n$.
\heading{Volume preservation}
Taking determinants gives $(\det M_t)^2=1$. Continuity in time and $\det M_0=1$
select $\det M_t=1$.
\medskip
Alternatively, $\operatorname{tr}(\mathbb J_nQ_t)=0$ because one factor is
skew-symmetric and the other symmetric. The Jacobian formula gives the same result.
\slidecite{Manuscript \S\,1.2.3, Assessment 10.4. \hyperlink{main:ham-conservation}{Return to geometric conservation.}}
\end{frame}

\begin{frame}[t]{Appendix A13: the entropy rate}
\hypertarget{app:entropy}{}
Under the smoothness, finiteness, and boundary assumptions of the entropy theorem,
\begin{equation*}
  \begin{aligned}
    \frac{\dd}{\dd t}S(\rho_t)
    &=-\int(1+\log\rho_t)\partial_t\rho_t\dd x\\
    &=\int(1+\log\rho_t)\nabla\cdot(\rho_tv_t)\dd x\\
    &=-\int\rho_tv_t\cdot\nabla\log\rho_t\dd x\\
    &=-\int v_t\cdot\nabla\rho_t\dd x\\
    &=\int\rho_t\nabla\cdot v_t\dd x.
  \end{aligned}
\end{equation*}
\heading{Interpretation along the flow}
The density formula gives $\log\rho_t(x_t)=\log\rho_0(x_0)-\log J_{0,t}(x_0)$.
When these terms are integrable,
\begin{equation*}
  S(\rho_t)-S(\rho_0)=\mathbb E_{\rho_0}[\log J_{0,t}].
\end{equation*}
\slidecite{Manuscript \S\,1.2.4, Assessment 9.10. \hyperlink{main:entropy}{Return to differential entropy.}}
\end{frame}

\begin{frame}[t]{Appendix A14: the Gibbs variational identity}
\hypertarget{app:gibbs}{}
For bounded $U$, the laws $\nu_U=Z_U^{-1}e^{-U}\nu$ and $\nu$ are equivalent.
If $\mu\ll\nu$, the Radon--Nikodym chain rule gives
\begin{equation*}
  \frac{\mathrm d\mu}{\mathrm d\nu_U}
  =Z_Ue^U\frac{\mathrm d\mu}{\mathrm d\nu},\qquad
  \log\frac{\mathrm d\mu}{\mathrm d\nu_U}
  =\log\frac{\mathrm d\mu}{\mathrm d\nu}+U+\log Z_U.
\end{equation*}
Integrating under $\mu$ yields
\begin{equation*}
  \KL(\mu\Vert\nu_U)=\KL(\mu\Vert\nu)+\int U\dd\mu+\log Z_U.
\end{equation*}
If $\mu\not\ll\nu$, both relative entropies are infinite. Bounded $U$ keeps the
identity meaningful in the extended-real sense.
\heading{Minimization}
\begin{equation*}
  \KL(\mu\Vert\nu)+\int U\dd\mu\geq-\log Z_U,
\end{equation*}
with equality precisely when $\KL(\mu\Vert\nu_U)=0$, or $\mu=\nu_U$.
\slidecite{Manuscript \S\,1.2.4, Assessment 11.5. \hyperlink{main:gibbs}{Return to Gibbs tilting.}}
\end{frame}
